\chapter{Breathing Practices and Pranayama}
\label{ch:pranayama}

\emph{How does a beginner learn to breathe well, and which pranayamas must the candidate perform?}
The syllabus puts two layers under this heading. First come the preparatory breathing practices, which make the breath slow, deep and smooth. Then come three pranayamas done without holding the breath: nadi shodhana (also called anuloma viloma), sheetali and bhramari. This chapter teaches both layers and the words a candidate must know. The meaning and place of pranayama in yoga are in the introduction to pranayama and dhyana (see \cref{ch:pranayama-intro}).

\section{Preparatory breathing}

\subsection*{Why practise first}
In a normal breath the in-breath is active and the out-breath passive and relaxing. Most people under stress breathe quickly and high in the chest. Before pranayama, the student learns to move each part of the trunk with the breath, so that the breath becomes slow, deep and smooth. The Board syllabus lists the abdominal, thoracic and clavicular breaths and the full yogic breath under pranayama, and the guidebook calls the full yogic breath a basic step in restoring natural breathing.

\subsection*{Position and rules}
Lie on the back, or sit on a chair or cross-legged, with the head, neck and spine in one line. Close the eyes, relax, and breathe through the nose, not the mouth. Keep the breath slow and comfortable, never forced. Stay relaxed in body and mind.

\subsection*{Abdominal (diaphragmatic) breathing}
Place the palms on the belly, the middle fingers near the navel. Breathe in so that the diaphragm moves down and the belly swells without pressure; breathe out slowly so that the belly sinks. The chest stays still. A common fault is to draw the belly in while breathing in; it should rise. This is the slow, relaxing breath the guidebook recommends for breaking the stress cycle.
\ycbfig{ycb1-fig-abdominalbreathing}

\subsection*{Thoracic (chest) breathing}
Place the palms lightly on the sides of the ribs. Breathe in so that the ribs widen in all directions; breathe out so that they come back. The belly stays still and the shoulders stay relaxed. It expands and contracts the rib cage and uses the middle part of the lungs. The guidebook notes that it takes in less air than the abdominal breath and tends to become quick, and that it takes over under stress.
\ycbfig{ycb1-fig-chestbreathing}

\subsection*{Clavicular (collarbone) breathing}
Place the fingertips on the collarbones. Breathe in so that the upper chest and collarbones rise a little; breathe out so that they lower. The shoulders do not move. Only the top of the lungs fills, so the breath is small and shallow. It is typical of stress or difficult breathing, so it is practised only as a step towards the full breath, never on its own.
\ycbfig{ycb1-fig-clavicularbreathing}

\subsection*{The full yogic breath}
The three breaths join into one smooth wave, also called three-part or combination breathing. Breathe in through the belly, then the ribs, then the collarbones, without a break; breathe out in the reverse order, slowly. The parts are no longer separate. Hold the breath out for a few seconds at the end of the out-breath; that completes one round. Begin with 5 to 10 rounds and build to 10 minutes a day. Remember \emph{belly, ribs, collarbones} on the way in and the reverse on the way out.

\begin{caution}
The guidebook says that light-headedness after the full breath shows it is done right. Teach the opposite: if a student feels dizzy, stop and let the breath return to normal. Keep any pause short, and only if comfortable.
\end{caution}
\ycbfig{ycb1-fig-threepartbreathing}

\section{Pranayama}

\subsection*{The word and the phases}
The guidebook splits \emph{pranayama} into \emph{prana}, the life force, and \emph{ayama}, expansion or control, and reads it as control of the life force through the breath. It is the fourth of Patanjali's eight limbs, after asana and before pratyahara. Patanjali defines it as the breaking of the flow of inhalation and exhalation (\emph{tasmin sati śvāsa-praśvāsayor gati-vicchedaḥ prāṇāyāmaḥ}).

Three words describe the phases of a breath:
\begin{itemize}
\item \textbf{Puraka}: the filling phase, the in-breath.
\item \textbf{Kumbhaka}: the held phase. The guidebook names two types, sahita and kevala.
\item \textbf{Rechaka}: the emptying phase, the out-breath.
\end{itemize}
The three pranayamas of the Level 1 syllabus are done without kumbhaka.

\subsection*{Nadi shodhana or anuloma viloma (alternate nostril breathing)}
\begin{enumerate}
\item Sit in a meditative posture with the spine and head upright and the eyes closed. Relax with a few deep breaths.
\item Rest the left hand on the left knee in jnana mudra.
\item Hold the right hand in nasagra mudra: the index and middle fingers folded, the thumb by the right nostril, the ring and little fingers by the left nostril.
\item Breathe in through the left nostril. Close the left, open the right, and breathe out through the right.
\item Breathe in through the right. Close the right, open the left, and breathe out through the left. That is one round.
\end{enumerate}
The protocol gives five rounds. Beginners keep the in-breath and the out-breath equal; later the out-breath is made twice as long. The breath stays slow, steady and unforced. Slides that suggest five to ten minutes a day go beyond the protocol; use the protocol count in the exam. Nadi shodhana is traditionally credited with a steady mind and lower stress.
\ycbfig{ycb1-fig-anulomvilom}

\subsection*{Sheetali (cooling breath)}
Sit in a meditative posture. Roll the tongue into a tube and let it rest just outside the lips. Breathe in slowly through the tube; draw the tongue in, close the mouth, and breathe out slowly through the nose. Never force the breath, and stop if dizzy. The word means cooling, and the practice is credited with cooling and calming body and mind. If the tongue cannot be rolled, use sheetkari: breathe in through the teeth with a soft hissing sound.

\begin{caution}
Sheetali can lower the blood pressure, so people with low blood pressure avoid it. Avoid it also with a cold, cough, asthma or tonsillitis, and in cold weather.
\end{caution}
\ycbfig{ycb1-fig-sitali}

\subsection*{Bhramari (humming bee breath)}
Sit in a meditative posture with the eyes closed. Breathe in deeply through the nose, then breathe out slowly and steadily with a deep humming sound like a bee. That is one round; the protocol gives five. In a second form, after the in-breath the index fingers rest on the closed eyes, the ring and little fingers close the lips, and the thumbs close the ears (shanmukhi mudra); the humming out-breath follows. Bhramari is credited with calming the mind, easing stress and anxiety, and preparing for concentration and meditation.
\textit{Caution:} avoid it with a nose or ear infection.
\ycbfig{ycb1-fig-bhramari}

\subsection*{Other pranayamas to know by name}
\begin{itemize}
\item \textbf{Ujjayi (victorious breath)}: breathe in and out with a gentle narrowing of the throat, which makes a soft, even, hissing sound. It can be done sitting, standing or lying, and is used during asanas to help focus.
\item \textbf{The eight Hatha Pradipika kumbhakas}: surya bhedana, ujjayi, sitkari, sitali, bhastrika, bhramari, murchha and plavini. By puraka, rechaka and kumbhaka there are three kinds, but as accompanied or alone there are two, sahita and kevala.
\item \textbf{The eight of the Gheranda Samhita}: sahita, surya bheda, ujjayi, sitali, bhastrika, bhramari, murchha and kevali.
\item \textbf{Patanjali's four kinds}: external, internal and suspended, each regulated by place, time and number; and a fourth that goes beyond the external and internal.
\item \textbf{Cooling and warming}: sheetali and sheetkari cool, chandra bhedana (through the left nostril) cools, and surya bhedana (through the right) warms. No text classes bhastrika as a heating pranayama; the Hatha Pradipika says it increases the digestive fire, a traditional claim. The guidebook advises people with heart disease or high blood pressure to avoid forceful bhastrika.
\item \textbf{Plavini and sitali}: in plavini the belly is filled with air so that the body floats on water like a lotus leaf. The Hatha Pradipika credits sitali, not plavini, with removing hunger and thirst. Both are traditional claims, and plavini is outside the syllabus.
\item \textbf{Kapalabhati} is a cleansing practice that works through the breath (see \cref{ch:cleansing}); it is not in the syllabus list of pranayamas.
\end{itemize}
\ycbfig{ycb1-fig-ujjayi}
\ycbfig{ycb1-fig-khechari}

\subsection*{Teaching pranayama}
Teach the hand position away from the face first, then add the breath. Keep every beginner practice short and comfortable; nothing here asks a student to strain or hold the breath. Sit a student on a chair or against a wall if the spine will not stay upright on the floor.

\begin{recall}
\begin{itemize}
\item Sectional breaths: abdominal, thoracic, clavicular; full yogic breath joins them, belly to collarbones.
\item Puraka in, kumbhaka held, rechaka out; the syllabus practices are without kumbhaka.
\item Nadi shodhana: start in through the left; five rounds; equal breath for beginners.
\item Sheetali: tongue rolled, in through the tongue, out through the nose; avoid with low blood pressure.
\item Bhramari: humming out-breath; five rounds; second form with shanmukhi mudra.
\end{itemize}
\end{recall}

\subsection*{Worked examples}
\begin{itemize}
\item \textbf{A student counts the breath and becomes tense.} Drop the counts and keep the breath equal and gentle for several rounds. Return to counting only when the student asks for it.
\item \textbf{A student feels dizzy during the full yogic breath.} Stop, let the breath return to normal, and resume later with smaller, slower breaths. Dizziness is a sign to stop, not a sign of success.
\item \textbf{It is a cold morning and a student has a cough.} Leave out sheetali; give bhramari or nadi shodhana instead.
\end{itemize}

\subsection*{Check yourself}
\begin{enumerate}
\item Name the three sectional breaths.
\item In which order do the sections fill in the full yogic breath?
\item What is another name for abdominal breathing?
\item Why are the sectional breaths practised before pranayama?
\item Name the three phases of a breath in pranayama.
\item In nadi shodhana, which nostril takes the first in-breath?
\item How many rounds does the protocol give for nadi shodhana and bhramari?
\item What ratio does a beginner use in nadi shodhana?
\item Give two cautions for sheetali.
\item What sound does ujjayi make?
\end{enumerate}

\subsection*{Answers}
\begin{enumerate}
\item Abdominal, thoracic (chest) and clavicular.
\item Belly, chest, collarbones; the out-breath is the reverse.
\item Diaphragmatic breathing.
\item To make the breath slow, deep and smooth.
\item Puraka (in), kumbhaka (held), rechaka (out).
\item The left.
\item Five rounds each.
\item In-breath and out-breath of equal length.
\item Any two of: low blood pressure, cold, cough, asthma, tonsillitis, cold weather.
\item A soft, even hissing sound made by gently narrowing the throat.
\end{enumerate}

\paragraph{Sources.} Common Yoga Protocol, Ministry of AYUSH (nadi shodhana, sheetali, bhramari, rounds); the official guidebook for Level I (2016, older scheme), chapter on pranayama and meditation (sectional breathing, kumbhaka types, ujjayi); Swami Satyananda Saraswati, \emph{Asana Pranayama Mudra Bandha} (Bihar School of Yoga): sectional and yogic breathing; Hatha Pradipika (Pancham Sinh) and Gheranda Samhita (Srisa Chandra Vasu), The Sacred Books of the Hindus vol.~XV (1915): the named lists, bhastrika, plavini and sitali; Yoga Sutra 2.50 and 2.51 (J.~H. Woods translation): the four kinds; the Board syllabus for Level 1.
